\section{Alternative coverage and area calculations}
\label{ext:alternate-estimates}

The fixed-margin criterion already proves coverage for the cap choices
in Section~\ref{ext:cap-section}.  The following direct calculations
explain three particular features: why squared-radius caps suit a
comparison of radial squares, how a Cartesian discriminant avoids
locating the perturbed interface, and how node-centred sectors retain
the complementary radial cancellation.  They are independent
explanations of those cases, rather than inputs to the main construction
or the vanishing-margin argument.

\subsection{Squared radii}
\label{ext:squared-direct}
Let $n$ be a positive integer, put $k=1/n$, use $n^2$ equal intervals,
and write their width as
$\Delta=2k^2$.  Put
\[
 M_I=(l+r)/2=2\beta_I,\qquad B_I=(1+lr)/2=2\alpha_I,
 \qquad f(q)=(q^2+q^4)/4.
\]
The squared-radius cap \eqref{ext:squared-cap}, with
$\tau=2k$ and $\eta=1/960$, is
$R(q)^2=1/4+k^2(f(q)+1/240)$.
The elementary secant identity is exact:
\begin{equation}\label{ext:squared-identity}
 2B_I-M_I^2=1-\frac{\Delta^2}{4}.
\end{equation}
For the two angularly selected intercepts, write
\[
 S=s-kxM_I,\qquad y=qS+kxB_I,
 \qquad T=1-s+kyM_J,\qquad x=pT-kyB_J.
\]
The crossing argument supplies $S,T\ge0$.  If both radial caps fail,
then $S,T>1/2$ and $u=s-1/2=O(k)$.
Since $M_I=q+O(k^2)$ and $M_J=p+O(k^2)$, substitution in
\[
 S^2=s^2-2kxSM_I-k^2x^2M_I^2,
 \qquad
 T^2=(1-s)^2+2kyTM_J-k^2y^2M_J^2
\]
and use of \eqref{ext:squared-identity} give the uniform expansions
\begin{align}
 S^2&=\tfrac14+u+u^2-2kxy+k^2x^2+O(k^3),\nonumber\\
 T^2&=\tfrac14-u+u^2+2kxy+k^2y^2+O(k^3).
 \label{ext:squared-expansions}
\end{align}
For example $SM_I=y-kxB_I+O(k^2)$, and the second identity uses
$TM_J=x+kyB_J+O(k^2)$.
Since $|M_I|,|M_J|\le1$, the failures $S,T>1/2$ first give
$|u|\le k$.  The first strict lower bound on $S^2$ then gives
$u-2kxy\ge-Ck^2$, and the second gives
$u-2kxy\le Ck^2$, for a uniform constant $C$.
Thus $u=2kxy+O(k^2)$.  Adding \eqref{ext:squared-expansions} now yields
\[
 \frac{S^2+T^2-1/2}{k^2}
 =x^2+y^2+8x^2y^2+O(k).
\]
On the other hand $q=2y+O(k)$ and $p=2x+O(k)$, so
\[
 f(q)+f(p)=x^2+y^2+4(x^4+y^4)+O(k)
 \ge x^2+y^2+8x^2y^2+O(k).
\]
The failed caps require the previous quotient to exceed
$f(q)+f(p)+2/240$, a contradiction for small $k$.

Let $C_k$ denote the sum of the perpendicular base areas of all $2n^2$
cylinders. The two families have equal total area, and the square in the
cap makes the radial integration exact before the projection factor is
expanded:
\[
 \frac{C_k}{2H}=
 \sum_I\int_I
 \frac{\frac12[\frac14+k^2(f(q)+1/240)]}
 {\sqrt{1+k^2(B_I^2+2M_I^2)}}\,dq
 =\frac14+k^2\left(\frac1{240}-\frac1{120}\right)+O(k^4).
\]
The squared cap therefore supplies both the coverage comparison and
the exact radial part of the area integral.

\subsection{A Cartesian discriminant}
\label{ext:cartesian-direct}
For an equal sector $I=[l,r]$ of width $\Delta$, the coefficients
\eqref{ext:ordinary-secants} satisfy
\[
 \beta_I^2-\alpha_I+1/4=\Delta^2/16.
\]
Suppose the angular crossing for one family has given
\[
 S=T-\delta\beta_I\ge0,\qquad
 Y_0=Y-\delta\alpha_I,\qquad lS\le Y_0\le rS.
\]
Here $(T,Y,\delta)=(s,y,\tau x)$ in the first family and
$(1-s,x,-\tau y)$ in the second.  Because $2\beta_I$ is the sector
midpoint, $|2\beta_IS-Y_0|\le\Delta S/2$.  Consequently
\begin{align}
 Q&:=T^2-\delta Y+\delta^2/4\nonumber\\
  &=S^2+\delta(2\beta_IS-Y_0)+\delta^2\Delta^2/16
    \ge (S-|\delta|\Delta/4)^2.
 \label{ext:cartesian-discriminant}
\end{align}
In particular $\sqrt Q\ge S-|\delta|\Delta/4$.
Take the Cartesian cap \eqref{ext:cartesian-aperture} with $b=1$,
fixed $\eta>0$, and $\Delta=O(\tau^2)$.
Since $0\le\alpha_I\le1/2$, $|Y|\le1$ and $|\delta|\le\tau$,
we have $|Y_0|\le1+\tau/2$.  For $\tau\le1$, this gives
$D(Y_0)\le D(3/2)=45/8$.  Hence the polynomial upper bound is less
than one whenever $\tau^2(45/8+\eta)<1/2$.
Failure of the selected cylinder therefore
implies $S>1/2+\tau^2(D(Y_0)+\eta)$ even if the outer bound $S\le1$
failed.  Since $Y_0=Y+O(\tau)$ and $|\delta|\Delta=O(\tau^3)$,
\eqref{ext:cartesian-discriminant} gives
\[
 \sqrt Q>\tfrac12+\tau^2\bigl(D(Y)+\eta/2\bigr)
\]
for sufficiently small $\tau$, uniformly in both families.
If both failed, squaring these positive lower bounds and discarding
their nonnegative fourth-order terms would give
\begin{align*}
 s^2&>\tfrac14+\tau xy+
          \tau^2\bigl(D(y)+\eta/2-x^2/4\bigr),\\
 (1-s)^2&>\tfrac14-\tau xy+
          \tau^2\bigl(D(x)+\eta/2-y^2/4\bigr).
\end{align*}
The right sides are positive for small $\tau$.  Expanding their square
roots with $\sqrt{1/4+v}=1/2+v-v^2+O(v^3)$ shows that their sum is
\[
 1+\tau^2\bigl(D(x)+D(y)+\eta-(x^2+y^2)/4-2x^2y^2\bigr)
       +O(\tau^3)
 =1+\tau^2\bigl((x^2-y^2)^2+\eta\bigr)+O(\tau^3)>1,
\]
contradicting $s+(1-s)=1$.

For either outer bound $b\in\{3/4,1\}$, the Cartesian aperture admits a
direct area calculation without changing to angular variables at the
boundary.  In the intercept coordinates $(Y,S)$, put
\[
 \mathcal D_\tau=\{0\le S\le b,\ |Y|\le S,
            \ S\le1/2+\tau^2(D(Y)+\eta)\},\qquad
 \mathcal D_0=\{0\le S\le1/2,\ |Y|\le S\}.
\]
For $|Y|\le1/2$ the added vertical length is exactly
$\tau^2(D(Y)+\eta)$ for small $\tau$.  Each remaining side region has
both width and height $O(\tau^2)$, because
$1/2<|Y|\le S\le1/2+O(\tau^2)$ there.  Thus
\[
 |\mathcal D_\tau|=\frac14+
 \tau^2\int_{-1/2}^{1/2}(D(Y)+\eta)\,dY+O(\tau^4).
\]
Expanding the projection factor over $\mathcal D_0$, and using
$Y=qS$ in its quadratic coefficient, gives the one-family cost
\[
 \frac14+\tau^2\left[
 \underbrace{\int_{-1/2}^{1/2}D(Y)\,dY}_{1/30}+\eta
 -\underbrace{\frac12\int_0^{1/2}S\,dS
   \int_{-1}^1\left(\frac{(1+q^2)^2}{16}+\frac{q^2}{2}\right)dq}_{17/480}
 \right]+O(\tau^4).
\]
Here the cost is divided by $H$ in the edge-two normalization.
The replacement of $\mathcal D_\tau$ by $\mathcal D_0$ changes the
quadratic coefficient by $O(\tau^2)$, and the coefficient sampling
error is also $O(\tau^2)$; both contribute only $O(\tau^4)$ to the
area.  Thus the normalized two-family cost is
$1/2+(2\eta-1/240)\tau^2+O(\tau^4)$, as in
\eqref{ext:general-cost}.  The Cartesian entries in
Appendix~\ref{ext:scaling-specializations} follow by substitution; the
edge-one discriminant and aperture are translated there as well.

\subsection{A direct check for node-centred sectors}
\label{ext:endpoint-direct}

Use the node-centred sectors of Section~\ref{ext:endpoint-sectors},
and write their tilt as $t=1/k$.  The radial cutoff at a sample node $q$
is $1/2+t^2(d(q)+1/1000)$, and each angular point of its sector lies
within $t^2$ of $q$.  If the two selected sample nodes are $q,p$, write
\[
 S=s-txq/2,\quad y=qS+tx\alpha(q)+e,
 \qquad
 T=1-s+typ/2,\quad x=pT-ty\alpha(p)+e_*,
\]
where $|e|\le St^2$, $|e_*|\le Tt^2$, and
$\alpha(q)=(1+q^2)/4$.  Direct substitution gives
\begin{align*}
 s-txy&=\tfrac12+(S-\tfrac12)(1-tqx)-t^2\alpha(q)x^2-txe,\\
 1-s+txy&=\tfrac12+(T-\tfrac12)(1+tpy)-t^2\alpha(p)y^2+tye_*.
\end{align*}
If both radial caps failed along a sequence $t\to0$, then $S,T>1/2$ and
$S+T=1+t(yp-xq)/2=1+O(t)$, so $S,T\to1/2$.
Passing to a subsequence with $q\to q_0$, $p\to p_0$ gives
$y\to q_0/2$, $x\to p_0/2$, while $e/t,e_*/t\to0$.
Adding the identities and dividing by $t^2$ would therefore imply
\[
 0\ge\frac2{1000}+\frac{(q_0^2-p_0^2)^2}{16}>0.
\]
This complementary identity gives a second explanation for why the
half-width endpoint sectors still cover the boundary.
